%% kapa: exam=1 round=33 year=2022 period=1 form=A subject=civil_law subject_ko=민법 seq=22 number=22
%% source: Q-Net 감정평가사 기출문제 (cst00309) artlSeq=5209997 file=raw/1008/5209997/2022년도 제33회 감정평가사 1차시험 1교시(민법,경제학,부동산학).pdf
%% conversion: kordoc 4.24.51 -> tools/gen1.py (auto); review=proofread-sonnet; flags=none
\begin{question}{22}
물권적 청구권에 관한 설명으로 옳지 않은 것은? (다툼이 있으면 판례에 따름)
\begin{choices}
  \item 지역권자는 지역권을 방해하는 자에 대하여 방해의 제거를 청구할 수 있다.
  \item 간접점유자는 제3자의 점유침해에 대하여 물권적 청구권을 행사할 수 있다.
  \item 직접점유자가 임의로 점유를 타인에게 양도한 경우에는 그 점유이전이 간접점유자의 의사에 반하더라도 간접점유자의 점유가 침탈된 경우에 해당하지 않는다.
  \item 부동산 양도담보의 피담보채무가 전부 변제되었음을 이유로 양도담보권설정자가 행사하는 소유권이전등기말소청구권은 소멸시효에 걸린다.
  \item 민법 제205조 제2항이 정한 점유물방해제거청구권의 행사를 위한 ‘1년의 제척기간’은 출소기간이다.
\end{choices}
\end{question}
\answer{4}
